\subsection{Data Contract v2}

La etapa de datos se cerró mediante un contrato ejecutable, no mediante una inspección manual
única. Los artefactos centrales son \artifact{configs/data_contract_v2.yaml} y
\artifact{tools/audit_data_contract_v2.py}. El audit completo produjo:
\[
46\ \text{PASS},\qquad
0\ \text{WARN},\qquad
0\ \text{FAIL},\qquad
1\ \text{SKIP}.
\]

Entre las invariantes comprobadas están:
\begin{itemize}
\item 197 IDs oficiales únicos;
\item partición exacta 138/59;
\item target visible sólo en entrenamiento;
\item formas BASIC $107{,}666\times96$ y PRO $47{,}804\times20$;
\item cobertura de las 197 parcelas en ambas tablas satelitales;
\item parsing estricto de fechas y unicidad parcela--fecha--sensor;
\item validez geométrica y CRS;
\item continuidad de rasters climáticos;
\item continuidad diaria CHIRPS;
\item metadatos/unidades WaPOR;
\item inventario SoilGrids y reglas de CRS;
\item cobertura SIAP 2003--2025, schemas y duplicados;
\item inventario CEM;
\item exclusión explícita de ERA5-Land como fuente no certificada.
\end{itemize}

El contrato funciona como \emph{gate}: si cambian fuentes crudas, primero debe volver a
ejecutarse el audit; ningún checkpoint posterior debería asumir que una descarga nueva conserva
las mismas unidades o cobertura.

\subsection{Inmutabilidad de fuentes y transformación versionada}

El repositorio separa:
\begin{enumerate}
\item evidencia fuente inmutable bajo \artifact{data/source/};
\item datos externos reproducibles/locales bajo \artifact{data/raw/external/};
\item artefactos derivados pequeños y versionados;
\item reportes generados por checkpoints.
\end{enumerate}

Las correcciones no se escriben sobre el archivo fuente. Por ejemplo, metadata defectuosa de
un NetCDF puede normalizarse en un artefacto derivado, pero el input original permanece
intacto. Esto permite reconstruir la cadena
\[
\text{source}
\longrightarrow
\text{loader}
\longrightarrow
\text{normalización}
\longrightarrow
\text{feature}
\longrightarrow
\text{modelo}.
\]

La separación tiene una consecuencia estadística: cualquier resultado puede auditarse contra
la versión exacta de la tabla que lo generó, reduciendo el riesgo de que un mismo nombre de
feature represente transformaciones distintas en checkpoints diferentes.

\subsection{Integration fixtures}

Se creó un fixture real de 12 parcelas con IDs compartidos entre fuentes. La selección incluye
dos parcelas por estrato Estado$\times$CONJUNTO, diversidad target-free y casos de borde
administrativos. El fixture contiene
\[
12\ \text{parcelas},\qquad
192\ \text{filas BASIC},\qquad
96\ \text{filas PRO},\qquad
201\ \text{filas SIAP},
\]
además de geometrías, clima, CHIRPS, topografía, SoilGrids, WaPOR y CEM. El conjunto ocupa
aproximadamente 3.77 MiB y permite probar joins/extractores sin requerir el archivo externo
completo.

La selección del fixture no usa \artifact{RENDIMIENTO_T_HA}. Esto es relevante porque incluso
un dataset de integración podría introducir un sesgo sutil si se eligieran ejemplos por
valores extremos de $y$.

\subsection{CRS y operaciones espaciales}

Las operaciones espaciales se realizan con una política explícita de CRS. EPSG:4326 es
adecuado para almacenamiento/intercambio de coordenadas angulares, pero no para interpretar
distancias euclidianas en grados como kilómetros. Por ello, cualquier cálculo métrico debe
transformarse o usar una fórmula geodésica apropiada.

En etapas posteriores, la componente geográfica de las distancias transductivas se deriva de
centroides parcelarios y se normaliza antes de combinarse con distancias agronómicas. La
consistencia de CRS es especialmente importante porque la geografía terminó siendo la señal
de transferencia más fuerte del proyecto.

\subsection{Missingness como semántica y no sólo como problema numérico}

La política de datos distingue tres casos:
\begin{enumerate}
\item valor ausente porque la fuente no cubre una parcela/fecha;
\item variable no definida para un sensor;
\item valor eliminado por QA.
\end{enumerate}
No todos los NaN significan lo mismo. Durante la construcción se eliminan columnas
completamente nulas porque no pueden aportar información. Las columnas parcialmente faltantes
se conservan y la imputación aprendida se delega al pipeline de modelado. De esta forma, una
decisión de ingeniería no utiliza accidentalmente distribución de validación.

\subsection{Hashing, manifest y trazabilidad}

Los builds derivados registran hashes de inputs, commit de Git y manifest de features. Esta
trazabilidad permite responder tres preguntas que son particularmente importantes en una
competencia con múltiples fuentes:
\begin{enumerate}
\item ¿qué archivo crudo originó una variable?;
\item ¿qué versión de código produjo la tabla usada por un modelo?;
\item ¿qué features pertenecían realmente al conjunto clean o competition en ese commit?
\end{enumerate}

La respuesta no depende de memoria humana ni de nombres de notebooks. El estado reproducible
queda materializado en YAML, JSON y CSV versionados.

\subsection{Criterios de cierre del Checkpoint 01}

Checkpoint~01 no produjo todavía un modelo ni una tabla final $197\times p$. Su cierre
significó algo más fundamental: cada fuente relevante tenía una semántica, una política de
unidades, una llave de unión, un tratamiento de missingness y un mecanismo de auditoría
documentados. El Checkpoint~02 pudo entonces operar sobre una base reproducible en vez de
incorporar limpieza implícita dentro del código de modelado.

\begin{decision}
Una feature sólo puede entrar a modelado si puede trazarse hasta una fuente y una regla de
transformación. La ausencia de provenance se trata como un defecto de datos, no como un
detalle de documentación.
\end{decision}
